\documentclass[11pt]{article}
\usepackage{amsmath,amssymb,amsthm}
\usepackage{hyperref}

\title{Bridge: Crystal Bases $\longleftrightarrow$ Cylindric Partitions\\[4pt]
\large via Quantum Affine Representation Theory}
\author{Notes (from structural hole analysis)}
\date{}

\begin{document}
\maketitle

\noindent
These three papers form a bridge between the \textbf{crystal bases}
and \textbf{cylindric partitions} clusters in the INSTINCT corpus.
They were identified by the citation-aware structural hole detector:
both Schilling and Korff cite Chari--Pressley, but from opposite
mathematical directions.  The shared object is the
\emph{Kirillov--Reshetikhin module} for quantum affine algebras.

\bigskip

%% ===================================================================
\section{Chari--Pressley: Twisted Quantum Affine Algebras (1996)}
%% ===================================================================

\textbf{Reference:} V.~Chari and A.~Pressley,
\emph{Twisted Quantum Affine Algebras},
\href{https://arxiv.org/abs/q-alg/9611002}{arXiv:q-alg/9611002} (1996).

\medskip
\noindent\textbf{Summary:}
Classifies the irreducible finite-dimensional representations of
\emph{twisted} quantum affine algebras $U_q(\hat{\mathfrak{g}})$
(types $A_{2n}^{(2)}$, $D_{n+1}^{(2)}$, etc.).  This classification
provides the algebraic foundation that both papers below build on.

\medskip
\noindent\textbf{Bridge role:}
The Chari--Pressley classification is the dictionary that makes it
possible to study the \emph{same} representations from two different
combinatorial perspectives:
\begin{itemize}
  \item Crystal bases (Schilling): encode representation-theoretic
    data as colored directed graphs.
  \item Cylindric partitions (Korff): encode structure constants of
    the fusion/Verlinde algebra as generating functions over partitions.
\end{itemize}

%% ===================================================================
\section{Schilling: Crystal Bases Side (2007)}
%% ===================================================================

\textbf{Reference:} A.~Schilling,
\emph{Combinatorial structure of Kirillov--Reshetikhin crystals of
type $D_n^{(1)}$, $B_n^{(1)}$, $A_{2n-1}^{(2)}$},
\href{https://arxiv.org/abs/0704.2046}{arXiv:0704.2046} (2007).

\medskip
\noindent\textbf{Summary:}
Constructs Kirillov--Reshetikhin (KR) crystals for twisted and
non-simply-laced types using $\pm$~diagrams and Dynkin diagram
automorphisms.  Verifies Sternberg's conjecture for these types.

\medskip
\noindent\textbf{Key concepts:}
\begin{itemize}
  \item \textbf{Kirillov--Reshetikhin crystals:} combinatorial models
    (crystal graphs) for the KR modules classified by Chari--Pressley.
    Each node is a tableau-like object; edges are crystal operators
    $\tilde{e}_i, \tilde{f}_i$.
  \item \textbf{$\pm$~diagrams:} a technique for constructing KR
    crystals of non-simply-laced type from those of simply-laced type
    via Dynkin diagram folding.
  \item \textbf{Demazure crystals:} subsets of affine crystals that
    correspond to Demazure modules; KR crystals are shown to be
    isomorphic to certain Demazure crystals.
\end{itemize}

%% ===================================================================
\section{Korff: Cylindric Partitions Side (2011)}
%% ===================================================================

\textbf{Reference:} C.~Korff,
\emph{Cylindric versions of specialised Macdonald functions and a
deformed Verlinde algebra},
\href{https://arxiv.org/abs/1110.6356}{arXiv:1110.6356} (2011).

\medskip
\noindent\textbf{Summary:}
Connects cylindric partition generating functions to the
representation theory of quantum $\widehat{\mathfrak{sl}}_n$
via the Yang--Baxter equation.  Shows that cylindric skew tableaux
compute structure constants in a deformed Verlinde algebra.

\medskip
\noindent\textbf{Key concepts:}
\begin{itemize}
  \item \textbf{Cylindric skew tableaux:} the combinatorial objects
    being enumerated.  Their generating functions give
    \emph{skew Macdonald functions} (cylindric generalisation).
  \item \textbf{Verlinde algebra:} the fusion ring of the
    WZW conformal field theory at level~$k$.  Its structure constants
    $N_{\lambda\mu}^\nu$ count cylindric tableaux of appropriate shape.
  \item \textbf{Yang--Baxter equation:} the integrability condition
    that makes the connection to quantum groups work.  Transfer
    matrices built from $R$-matrices of $U_q(\widehat{\mathfrak{sl}}_n)$
    generate cylindric partition functions.
  \item \textbf{Kirillov--Reshetikhin modules:} appear as the
    auxiliary spaces in the transfer matrix construction.  This is the
    direct link to Chari--Pressley and to Schilling's crystals.
\end{itemize}

%% ===================================================================
\section{The Bridge}
%% ===================================================================

The conceptual picture is:
\[
\boxed{\text{Crystal graphs}}
\xleftarrow{\text{Schilling}}
\boxed{\text{KR modules of } U_q(\hat{\mathfrak{g}})}
\xrightarrow{\text{Korff}}
\boxed{\text{Cylindric partitions}}
\]

\begin{itemize}
  \item \textbf{Left arrow:}
    Schilling constructs the \emph{crystal basis} of each KR module ---
    a purely combinatorial object (a coloured directed graph) that
    encodes the module's weight structure and tensor product
    decomposition.

  \item \textbf{Right arrow:}
    Korff shows that the \emph{transfer matrices} built from KR modules
    (via $R$-matrices and the Yang--Baxter equation) generate
    cylindric partition functions.  The trace over auxiliary KR modules
    produces generating functions for cylindric skew tableaux.

  \item \textbf{Centre:}
    Chari--Pressley's classification tells us which KR modules exist
    and how they decompose.  Without it, neither the crystal
    construction nor the transfer matrix approach would have a
    foundation.
\end{itemize}

\medskip
\noindent\textbf{Relevance to Conjecture~2.7:}
If the positivity of $Q_{n,c}(q)$ has a representation-theoretic
explanation, this bridge suggests where to look: the cylindric
partition generating functions are traces over KR modules, and the
crystal bases of those modules carry a natural grading by
the energy function.  The coefficients of $Q_{n,c}(q)$ might
count something in the crystal (e.g.\ highest weight elements
of a given energy).

\medskip
\noindent\textbf{Discovery context:}
This bridge was found by the structural hole detector
(\texttt{detect\_structural\_holes.py}), which identified
crystal-bases $\leftrightarrow$ cylindric-partitions as a
structural hole (concepts overlap but papers don't cite each other)
and flagged Chari--Pressley as the sole external paper cited by
both clusters.

\end{document}
